% Generated draft: requires mathematical and source review.
\documentclass[11pt]{article}
\usepackage[margin=25mm]{geometry}
\usepackage{fontspec}
\setmainfont{Latin Modern Roman}
\usepackage{amsmath,amssymb,mathtools,unicode-math}
\setmathfont{Latin Modern Math}
\usepackage{enumitem,needspace,xurl,hyperref}
\setlength{\parindent}{0pt}
\setlength{\parskip}{5pt}
\setlength{\emergencystretch}{3em}
\newcommand{\N}{\mathbb{N}}
\newcommand{\Z}{\mathbb{Z}}
\newcommand{\Q}{\mathbb{Q}}
\newcommand{\R}{\mathbb{R}}
\newcommand{\C}{\mathbb{C}}
\newcommand{\F}{\mathbb{F}}
\newcommand{\A}{\mathbb{A}}
\newcommand{\PP}{\mathbb{P}}
\newcommand{\E}{\mathbb{E}}
\newcommand{\eps}{\varepsilon}
\newcommand{\dd}{\,\mathrm{d}}
\begin{document}
\section*{UnsolvedMath ID: 700742}
\textbf{Source problem code:} catalogue entry
\subsection{Definitions and mathematical statement}
Let \(A\) and \(C\) be abelian groups, written additively. The source uses a scalar \(n\) without specifying its range; the formulas below interpret \(n\) as an integer, with its intended range left unresolved. For an abelian group \(G\), write
\[
G[n]=\{g\in G:ng=0\},\qquad nG=\{ng:g\in G\}.
\]
Assume \(A[n]=0\), meaning that \(na=0\) implies \(a=0\) for every \(a\in A\). Write \(\operatorname{Hom}(C,A)\) for the group of homomorphisms and \(\operatorname{Ext}(C,A)=\operatorname{Ext}^{1}_{\mathbb Z}(C,A)\) for the group of equivalence classes of abelian-group extensions
\[
0\longrightarrow A\longrightarrow X\longrightarrow C\longrightarrow 0.
\]
Here equivalence means an isomorphism of extensions inducing the identity on both \(A\) and \(C\).

The source supplies the exact sequence
\[
0\longrightarrow
\frac{\operatorname{Hom}(C,A)}{n\operatorname{Hom}(C,A)}
\xrightarrow{\rho}
\operatorname{Hom}(C/nC,A/nA)
\xrightarrow{\delta}
\operatorname{Ext}(C,A)[n]
\longrightarrow 0.
\]
The reduction map is
\[
\rho([\varphi])(c+nC)=\varphi(c)+nA.
\]
Since \(n(A/nA)=0\), precomposition with \(C\to C/nC\) gives the canonical identification
\[
\operatorname{Hom}(C/nC,A/nA)\cong\operatorname{Hom}(C,A/nA).
\]
For \(f\in\operatorname{Hom}(C,A/nA)\), let \(q:A\to A/nA\) be the quotient map and form the pull-back
\[
X_f=\{(a,c)\in A\times C:q(a)=f(c)\}.
\]
The corresponding extension is
\[
0\longrightarrow A\xrightarrow{i_f}X_f\xrightarrow{p_f}C\longrightarrow 0,
\qquad i_f(a)=(na,0),\quad p_f(a,c)=c.
\]
Thus \(\delta(f)\) is the class of this extension, using the displayed identification of Hom groups. The pull-back diagram is
\[
\begin{array}{ccccccccc}
0&\longrightarrow&A&\xrightarrow{i_f}&X_f&\xrightarrow{p_f}&C&\longrightarrow&0\\
&&\downarrow\scriptstyle{\mathrm{id}_A}&&\downarrow\scriptstyle{(a,c)\mapsto a}&&\downarrow\scriptstyle{f}&&\\
0&\longrightarrow&A&\xrightarrow{n}&A&\xrightarrow{q}&A/nA&\longrightarrow&0.
\end{array}
\]

\textbf{Requested classification.} Use this scheme to classify the extensions intended by the source. The source says only ``certain extensions of \(A\) by \(C\)'' and does not specify their subclass or the desired classification criterion. The displayed diagram fixes the orientation as kernel \(A\) and quotient \(C\); the classification request remains incomplete beyond that orientation.

The source singles out the case \(nC=0\) with \(A\) torsion-free, where torsion-free means that \(ma=0\), for \(m\in\mathbb Z\setminus\{0\}\), implies \(a=0\). It states that in this case
\[
\operatorname{Ext}(C,A)[n]=\operatorname{Ext}(C,A).
\]
This special case is described as interesting, without a more specific classification target.
\subsection{Short English statement}
For abelian groups \(A,C\) with \(A[n]=0\), the supplied exact sequence relates homomorphisms modulo \(n\) to extension classes annihilated by \(n\), represented by explicit pull-backs. A. Mader asks that this construction be used to classify certain extensions, particularly when \(nC=0\) and \(A\) is torsion-free. The intended subclass and classification criterion are not specified.
\subsection{Sources}
\begin{itemize}
\item Supplied local UnsolvedMath registry record, permanent ID \(700742\), exact source problem code \texttt{catalogue entry}, title \emph{Kourovka Notebook Problem 11.50}. The registry labels the problem \texttt{open}; this is supplied metadata, not an independently verified mathematical or literature assessment.
\item \emph{The Kourovka Notebook (Unsolved Problems in Group Theory)}, Problem 11.50, proposed by A. Mader, 11th issue (1990). The supplied background identifies the source edition as the 21st issue edition (2026), arXiv:1401.0300v48. The supplied collection description names E. I. Khukhro and V. D. Mazurov as editors.
\item E. L. Lady and A. Mader, \emph{J. Algebra} \textbf{140} (1991), 36--64, cited in the supplied statement. No article title, DOI, or URL is supplied.
\end{itemize}
\end{document}
